\chapter{参数化与非参数化模型}
\begin{quote}\small
\textit{本章摘要。}本章比较两种表达函数复杂度的方式：参数化模型用有限维参数概括规律，非参数化模型则通过核、邻域或函数分布保留更灵活的形状。高斯过程和核方法贯穿全文，用来说明先验结构、计算代价和小样本性能之间的权衡，并与上一章的贝叶斯推断自然衔接。
\end{quote}
上一章用参数后验解释有限证据下的未知规律，但一个分布再精确，也受它所依附的模型结构限制。本章进一步问：一定要先选有限个回归系数再描述函数吗？能否直接规定不同输入处函数值如何共同变化，再由观测决定局部形状？参数化与非参数化的区别，正落在这个结构选择上。

假设要根据转速预测温升，可以先找转速相近的历史记录取平均，这就是近邻方法的直观起点。核方法用相似程度决定记录之间怎样相互影响，高斯过程进一步描述可能的函数及其概率。方法越灵活，越需要考虑样本是否足够、计算是否承受得起。还要分清一个预测数值和一整个预测分布：两个方法给出相同的平均温升，并不意味着它们对预测有同样的把握。

\section{两个概念}
参数化模型用固定维度参数描述分布，例如线性回归$y=\bm{w}^{\mathsf T}x+b+\varepsilon$。无论样本数如何增加，参数维度仍由输入维度和模型设计决定。非参数化模型不预先固定同样意义上的有限参数维度，而允许模型复杂度随数据增长，例如$k$近邻、高斯过程和Dirichlet过程模型。

\paragraph{近邻回归：由局部常数假设得到估计。}
设$x_i\in\R^d$、$Y_i=f^*(x_i)+\varepsilon_i$，给定输入后噪声独立、均值为零且方差为$\sigma^2$。固定查询$x$，按预先规定距离选出$k$个近邻集合$N_k(x)$。若假设该邻域内函数近似常数$c$，最小化$\sum_{i\in N_k(x)}(Y_i-c)^2$的一阶条件为$2\sum_i(c-Y_i)=0$，所以
\begin{equation}
 \hat f_k(x)=\frac1k\sum_{i\in N_k(x)}Y_i.
\end{equation}
给定训练输入，条件偏差为$k^{-1}\sum_{i\in N_k(x)}f^*(x_i)-f^*(x)$，条件方差为$\sigma^2/k$。若$f^*$为$L_f$-Lipschitz函数，近邻最远距离为$r_k(x)$，则
\begin{equation}
 |\operatorname{Bias}(x)|\leq L_fr_k(x),\qquad
 \E[(\hat f_k(x)-f^*(x))^2\mid X]\leq L_f^2r_k(x)^2+\sigma^2/k.
\end{equation}
增大$k$降低噪声方差，却扩大邻域偏差。若局部密度为正且规则，典型半径量级为$(k/n)^{1/d}$，平衡两项得到$k$约为$n^{2/(d+2)}$、均方误差量级约为$n^{-2/(d+2)}$；这是基于光滑与密度条件的局部量级论证，而非对任何数据的保证。维度越高，邻域收缩越慢，这就是距离方法的维数困难。分类近邻则平均类别指示以估计局部类别概率。

一致性通常要求$k\to\infty$且$k/n\to0$，并附加分布与正则条件；固定$k$一般保留不可消失的方差。距离尺度、无关特征和跨工况混合都可能破坏局部相似假设。朴素查询需约$O(nd)$距离计算，快速索引在高维也可能失效。

\paragraph{高斯过程示例。}
\begin{figure}[htbp]
\centering
\begin{tikzpicture}[x=0.75cm,y=0.55cm]
 \draw[->] (0,0)--(8.0,0) node[right] {$x$};
 \draw[->] (0,-2.0)--(0,2.8) node[above] {$f(x)$};
 \draw[gray!45] plot[smooth] coordinates {(0,0.2)(1,0.35)(2,0.1)(3,-0.2)(4,0.0)(5,0.5)(6,0.85)(7,0.6)(8,0.9)};
 \draw[blue!70!black, thick] plot[smooth] coordinates {(0,0.0)(1,0.25)(2,0.55)(3,0.2)(4,-0.25)(5,0.1)(6,0.65)(7,0.35)(8,0.7)};
 \fill[black] (1,0.25) circle (1.5pt) (2,0.55) circle (1.5pt) (3,0.2) circle (1.5pt) (5,0.1) circle (1.5pt) (6,0.65) circle (1.5pt);
 \draw[blue!35, fill=blue!8, opacity=0.55] plot[smooth] coordinates {(0,-0.9)(1,-0.6)(2,-0.2)(3,-0.65)(4,-1.0)(5,-0.65)(6,-0.1)(7,-0.5)(8,-0.1)} -- plot[smooth] coordinates {(8,1.5)(7,1.35)(6,1.4)(5,1.1)(4,0.5)(3,1.0)(2,1.2)(1,1.1)(0,1.0)} -- cycle;
 \node[font=\scriptsize,anchor=west] at (5.8,2.0) {远离观测点，不确定性增大};
 \node[font=\scriptsize] at (3.5,-1.7) {观测点};
\end{tikzpicture}
\caption{高斯过程用观测点附近的函数相关性收缩后验方差；远离数据时预测区间通常变宽。}
\label{fig:kernel-gp-prediction}
\end{figure}
高斯过程（Gaussian process, GP）直接在函数空间上定义先验：任意有限输入集合对应的函数值服从联合高斯分布，记为$f\sim\mathcal{GP}(m,k)$。给定训练输入$X$和观测$y=f(X)+\varepsilon$，其中$\varepsilon\sim\mathcal{N}(0,\sigma^2I)$，新点$x_*$的后验均值为
\begin{equation}
 \mu_* = m(x_*)+k_*^{\mathsf T}(K+\sigma^2I)^{-1}(y-m(X)),\label{eq:kernel-gp-mean}
\end{equation}
后验方差为
\begin{equation}
 \sigma_*^2=k(x_*,x_*)-k_*^{\mathsf T}(K+\sigma^2I)^{-1}k_*.\label{eq:kernel-gp-variance}
\end{equation}
这里$X$的$n$行分别为训练输入的转置，$y\in\R^n$为观测列向量，$m$为先验均值函数，$m(X)=[m(x_1),\ldots,m(x_n)]^{\mathsf T}$，$K_{ij}=k(x_i,x_j)$，$k_*=[k(x_1,x_*),\ldots,k(x_n,x_*)]^{\mathsf T}$。噪声与潜在函数独立。式\eqref{eq:kernel-gp-variance}针对潜在函数值$f(x_*)$；若预测含独立同方差噪声的新观测，还需加上$\sigma^2$。核函数$k$表达“相似输入应有相似函数值”的先验。远离训练样本时，预测方差通常增大，这提供了自然的不确定性信号。

图\ref{fig:kernel-gp-prediction}将观测点、预测曲线与不确定性带放在同一坐标中，提示读者区分均值拟合和区间宽度；远处是否回到先验还取决于所选核的相关结构。
\paragraph{GP后验不是另一个需要记忆的公式。}
设$C=K+\sigma^2I_n$，训练中心化观测$r=y-m(X)$，新点中心化函数值$g_*=f(x_*)-m(x_*)$。GP定义和独立高斯噪声给出
\begin{equation}
 \begin{bmatrix}r\\g_*\end{bmatrix}\sim\mathcal N\left(0,
 \begin{bmatrix}C&k_*\\k_*^{\mathsf T}&k_{**}\end{bmatrix}\right),
 \qquad k_{**}=k(x_*,x_*).
\end{equation}
构造残差$u=g_*-k_*^{\mathsf T}C^{-1}r$，直接计算$\operatorname{Cov}(u,r)=k_*^{\mathsf T}-k_*^{\mathsf T}C^{-1}C=0$。联合高斯中不相关意味着独立，且
\begin{equation}
 \Var(u)=k_{**}-k_*^{\mathsf T}C^{-1}k_*.
\end{equation}
因此给定$r$时$u$仍为同一零均值高斯，恢复先验均值便得到式\eqref{eq:kernel-gp-mean}与式\eqref{eq:kernel-gp-variance}。这个证明同时说明减去的项为何非负：$C$正定时$k_*^{\mathsf T}C^{-1}k_*\geq0$；整个条件方差非负则由联合协方差的半正定性保证。

对$m$个新点，$K_{X*}\in\R^{n\times m}$、$K_{**}\in\R^{m\times m}$，后验协方差为$K_{**}-K_{X*}^{\mathsf T}C^{-1}K_{X*}$。各点之间通常相关，不能把多个边际区间当作整条曲线的同时置信带。若核超参数固定，此协方差只依赖输入而不依赖观测值；异常的$y$不会自动增大方差。

$\sigma^2>0$使半正定$K$对应的$C$正定；无噪声且有重复输入时需处理奇异性。只有当远处$k_*$趋近零时，后验才回到先验均值和先验方差；周期核等结构可能在距离很远时仍高度相关。因此“远离数据更不确定”是常见核下的行为，而非GP定义本身的普适结论。

\section{核方法、函数空间与计算边界}
支持向量机展示了核函数如何把线性判别推广到非线性空间，稀疏高斯过程则以诱导变量降低函数后验近似的计算成本\cite{cortes1995svm,titsias2009sparsegp}。核化与贝叶斯函数建模可以共享工具，但优化目标和不确定性含义不同。
\paragraph{联系而非对立。}
参数化和非参数化并不是“现代”和“传统”的划分。神经网络是有限参数模型，但宽网络和核方法之间存在理论联系；贝叶斯神经网络在有限参数空间上保留后验；核方法则把参数学习转化为函数空间中的推断。真正的选择依据是数据量、计算资源、先验结构、预测任务和对不确定性的要求。

\begin{remark}
 当数据量很少而风险可解释性重要时，结构清晰的参数化概率模型往往更合适；当关系复杂且数据充足时，深度参数化模型可以提供更强的表示能力。两者都必须通过验证数据和任务指标检验。
\end{remark}

\paragraph{参数模型中的先验结构。}
参数化并不意味着没有先验。线性模型假设输入与输出的关系可由有限维方向表达；多项式模型增加曲率；状态空间模型把时间依赖写入转移方程。参数越少通常越容易估计，但若结构假设错误，偏差会成为主要误差来源。选择参数化模型就是选择哪些变化被认为值得显式表达。

\paragraph{核方法的函数空间视角。}
核函数$k(x,x')$定义了输入相似性。核岭回归的预测可以写成训练样本核函数的线性组合：
\begin{equation}
 f(x)=\sum_{i=1}^{n}\alpha_i k(x_i,x).
\end{equation}
核技巧把高维特征映射隐含在内积中，但训练通常需要处理$n\times n$核矩阵，样本规模增大时计算和存储成为瓶颈。稀疏近似、随机特征和 inducing points 可以降低这一代价。

\paragraph{什么相似性可以成为核。}
合法实值核必须对称，并对任意有限输入和$a\in\R^n$满足$\sum_{ij}a_ia_jk(x_i,x_j)\geq0$，即Gram矩阵半正定。若$k(x,x')=\phi(x)^{\mathsf T}\phi(x')$，该二次式等于$\|\sum_ia_i\phi(x_i)\|^2$，因此自动满足条件；任意手写距离的负值却不一定是合法核。

平方指数核$k(x,x')=s_f^2\exp[-\frac12\sum_j(x_j-x'_j)^2/\ell_j^2]$用$\ell_j>0$控制各维相关长度，$s_f^2>0$控制函数幅度。长度很大意味着该维变化难以使函数改变，长度很小则允许快速波动。核的非负加权和仍为核，可组合平滑、线性或周期结构，但组合之前应确定各分量对应什么物理变化。

\paragraph{核岭回归的表示与系数求解。}
在核对应的再生核Hilbert空间$\mathcal H_k$中，再生性质为
\begin{equation}
 f(x_i)=\langle f,k(x_i,\cdot)\rangle.
\end{equation}
考虑$\lambda>0$时的目标
\begin{equation}
 J(f)=\frac1{2n}\sum_i(y_i-f(x_i))^2+\frac\lambda2\|f\|_{\mathcal H_k}^2.
\end{equation}
把$f$正交分解为训练核函数张成空间中的$f_\parallel$与$f_\perp$。再生性质给出$f_\perp(x_i)=0$，故去掉$f_\perp$不改变训练拟合，却将范数平方减少$\|f_\perp\|^2$。最优解必可写为$f(\cdot)=\sum_i\alpha_i k(x_i,\cdot)$，这就是此处表示结论的推导。

代入后$J(\alpha)=\|y-K\alpha\|^2/(2n)+\lambda\alpha^{\mathsf T}K\alpha/2$，梯度条件为$K[(K+n\lambda I)\alpha-y]=0$。可选稳定解
\begin{equation}
 \alpha=(K+n\lambda I_n)^{-1}y.\label{eq:kernel-ridge-coefficients}
\end{equation}
若$K$奇异，不能在梯度式中随意约去$K$；但式\eqref{eq:kernel-ridge-coefficients}的解仍有效，系数的零空间差异表示同一个函数。正则化保证函数解唯一。零均值GP在$\sigma^2=n\lambda$时与核岭回归均值相同，但核岭目标本身不定义观测噪声或预测方差，不能直接把两者的不确定性等同。

\paragraph{从有限参数高斯先验得到核。}
令$f(x)=\phi(x)^{\mathsf T}w$，$\phi(x)\in\R^p$、$w\sim\mathcal N(0,S)$。任意有限函数值均为高斯线性变换，协方差为$k(x,x')=\phi(x)^{\mathsf T}S\phi(x')$。令$\Phi\in\R^{n\times p}$，则参数后验给出的均值与函数空间表达相等：
\begin{equation}
 (S^{-1}+\sigma^{-2}\Phi^{\mathsf T}\Phi)^{-1}\sigma^{-2}\Phi^{\mathsf T}
 =S\Phi^{\mathsf T}(\Phi S\Phi^{\mathsf T}+\sigma^2I)^{-1}.
\end{equation}
在$S$正定时，左右两边同乘左侧精度矩阵即可验证恒等式。它说明有限秩核GP也可有固定有限参数表示；所谓非参数性主要针对可随样本扩展的无限秩函数类，而不是“出现核就绝无有限参数”。

\paragraph{核支持向量机的对偶来源。}
对$y_i\in\{-1,1\}$，软间隔分类在特征空间中求
\begin{equation}
 \min_{w,b,\xi}\ \tfrac12\|w\|^2+C\sum_i\xi_i,
 \quad y_i(\langle w,\phi(x_i)\rangle+b)\geq1-\xi_i,
 \quad\xi_i\geq0.
\end{equation}
给两组不等式加入乘子$\alpha_i,\mu_i\geq0$。对$w,b,\xi_i$的驻点条件分别为$w=\sum_i\alpha_iy_i\phi(x_i)$、$\sum_i\alpha_iy_i=0$、$C-\alpha_i-\mu_i=0$。代回拉格朗日函数得
\begin{equation}
 \max_{0\leq\alpha_i\leq C,\ \sum_i\alpha_iy_i=0}
 \sum_i\alpha_i-\frac12\sum_{ij}\alpha_i\alpha_jy_iy_jk(x_i,x_j).
\end{equation}
凸性和可行条件使对偶可用于求解；预测得分为$\sum_i\alpha_iy_ik(x_i,x)+b$。若$0<\alpha_i<C$，互补松弛给出对应样本位于间隔边界，可用于计算$b$。非零系数对应支持向量。间隔得分不是故障概率，仍须单独校准；核矩阵半正定则是凸优化性质的重要前提。

\paragraph{高斯过程的边际似然。}
高斯过程不仅给出预测均值，还能通过边际似然选择核参数和噪声尺度：
\begin{equation}
 \log p(y\mid X)=-\frac12(y-m(X))^{\mathsf T}K_y^{-1}(y-m(X))-\frac12\log|K_y|-\frac n2\log(2\pi).\label{eq:kernel-marginal-likelihood}
\end{equation}
其中$K_y=K+\sigma^2I$，$|K_y|$表示行列式；$m\equiv0$时得到零均值先验的简式。
第一项衡量拟合，第二项惩罚过于灵活的函数空间。这个 Occam 因子体现了贝叶斯模型选择中“拟合与复杂度”的统一，但矩阵分解成本通常为$O(n^3)$。

\paragraph{边际似然的积分、导数与数值实现。}
由$f_X\sim\mathcal N(m_X,K)$及$y=f_X+\varepsilon$，独立高斯相加后$y\sim\mathcal N(m_X,C)$，$C=K+\sigma^2I$；把这个$n$维高斯密度取对数就得到式\eqref{eq:kernel-marginal-likelihood}。这一步已经积分掉函数值，并不是将最优函数值代入似然。

设超参数为$\eta$，$r=y-m_X$、$a=C^{-1}r$。利用$d\log|C|=\operatorname{tr}(C^{-1}dC)$和$dC^{-1}=-C^{-1}(dC)C^{-1}$，逐项微分得到
\begin{equation}
 \frac{\partial\log p(y\mid X)}{\partial\eta}
 =\frac12\operatorname{tr}\left[(aa^{\mathsf T}-C^{-1})\frac{\partial C}{\partial\eta}\right]
  +\left(\frac{\partial m_X}{\partial\eta}\right)^{\mathsf T}a.
\end{equation}
核长度、幅度和噪声可据此优化，正尺度通常取对数参数化。目标不保证凸，应检查多起点和边界。将超参数估计后固定属于经验贝叶斯，预测区间未包含超参数不确定性；充分贝叶斯做法还要对其后验积分。

数值上分解$C=LL^{\mathsf T}$，解$Lv=r$、$L^{\mathsf T}a=v$，可计算$r^{\mathsf T}C^{-1}r=v^{\mathsf T}v$及$\log|C|=2\sum_i\log L_{ii}$。预测时解$Lu=k_*$，方差为$k_{**}-u^{\mathsf T}u$，无需存储逆矩阵。浮点误差可能出现极小负方差，需结合残差和条件数诊断；不能把显著负值简单截断后忽略错误。训练分解为$O(n^3)$时间、$O(n^2)$存储，单点方差通常需$O(n^2)$求解。

\paragraph{行列式项不是独立的万能复杂度标尺。}
若$C=U\operatorname{diag}(c_j)U^{\mathsf T}$，中心化边际负对数似然为$\frac12\sum_j[(u_j^{\mathsf T}r)^2/c_j+\log c_j]+\frac n2\log(2\pi)$。增大某方向方差会减小残差惩罚，却增加归一化代价，两项共同决定折中。不能脱离数据与核谱，把行列式项简单解释成参数越多就惩罚越大。

\paragraph{非参数模型的计算边界。}
非参数模型可以表达复杂结构，却不代表复杂度无限且没有代价。kNN 的查询成本随样本数增加；核方法受核矩阵限制；高斯过程的后验更新需要处理越来越大的协方差矩阵。实际部署常用近似推断、原型压缩、局部模型或分层采样。

\paragraph{诱导变量如何降低代价。}
选$m\ll n$个诱导输入$Z$，令$u=f(Z)\in\R^m$。为简化取零均值，$K_{uu}\in\R^{m\times m}$可逆，$K_{Xu}\in\R^{n\times m}$。高斯条件公式给出
\begin{equation}
 p(f_X\mid u)=\mathcal N(K_{Xu}K_{uu}^{-1}u,\ K-Q),\qquad
 Q=K_{Xu}K_{uu}^{-1}K_{uX}.
\end{equation}
这里$K-Q$是诱导变量未解释的条件协方差，不能只因为使用低秩表示就断言它为零。取近似联合$q(f_X,u)=p(f_X\mid u)q(u)$，ELBO化为$\E_q\log p(y\mid f_X)-\KL(q(u)\|p(u))$。对高斯噪声，先在$p(f_X\mid u)$下取期望，得到
\begin{equation}
 \E[\log p(y\mid f_X)\mid u]
 =\log\mathcal N(y\mid K_{Xu}K_{uu}^{-1}u,\sigma^2I)
 -\frac{\operatorname{tr}(K-Q)}{2\sigma^2}.
\end{equation}
再对自由$q(u)$优化，利用$\sup_q\{\E_q\log h(u)-\KL(q\|p)\}=\log\int h(u)p(u)\,du$，可得
\begin{equation}
 \mathcal L=\log\mathcal N(y\mid0,Q+\sigma^2I)
 -\frac{\operatorname{tr}(K-Q)}{2\sigma^2}.
\end{equation}
迹修正惩罚遗漏的函数变异，因此这不是随意用$Q$替换$K$的精确模型。借助低秩恒等式主要成本可降至$O(nm^2+m^3)$、存储$O(nm)$；诱导点过少或位置不当会遗漏局部故障结构。$m=n$、诱导输入等于训练输入且矩阵适定时，$Q=K$恢复精确训练证据。

\paragraph{随机特征的另一路近似。}
对适当的平稳正定核，可将核写为频率分布下$\E_\omega\cos(\omega^{\mathsf T}(x-x'))$（幅度单独处理）。采样$\omega_j$和$b_j\sim\operatorname{Uniform}(0,2\pi)$，设$\phi_j(x)=\sqrt{2/m}\cos(\omega_j^{\mathsf T}x+b_j)$，由相位积分和独立平均有$\E[\phi(x)^{\mathsf T}\phi(x')]=k(x,x')$。这把核问题变成$m$维线性问题，但有限$m$引入随机近似误差；适用的频率分布由核决定，不能对任意相似性随意采样。

\paragraph{参数化与非参数化的混合。}
这些方法也可以组合使用。例如先用神经网络把振动窗口压成低维特征，再用 GP 预测并给出模型内的不确定性。这里若把编码器固定，GP 的区间通常不会自动包含编码器本身的不确定性。另一些任务只需在冻结特征上重新训练线性头，或用状态空间先验约束序列预测。组合前要说明每部分负责什么、哪些误差仍未计入，不能因为接上概率模型就认为整个系统已经校准。

\paragraph{模型选择案例。}
小样本温度预测可以从线性回归、核岭回归和高斯过程开始，比较留出设备上的误差、区间覆盖率、训练时间和增量更新成本。若数据规模扩大且输入包含图像或长序列，再考虑深度模型。这个顺序能够先验证任务是否可预测，再决定是否需要更复杂的表示学习。

\paragraph{本章小结。}
参数化模型以明确结构换取估计效率和解释性，非参数化模型以更灵活的函数空间换取计算成本。二者不是互斥阵营，关键是让模型复杂度、数据量、先验知识和不确定性需求相互匹配。

从局部常数的近邻估计，到再生核空间的正则解，再到高斯联合分布的条件化，本章展示了三种共享样本信息的机制。核规定哪些输入相关，正则或先验规定允许多大变化，近似方法规定在预算内保留哪些变化；任何一个环节都可能成为偏差来源。相同预测均值不意味着相同概率解释，较宽函数空间也不意味着在未覆盖工况下必然诚实地表达不确定性。

到这里，我们已经从参数分布走到了函数分布。下一章把读数按时间排好：平均温度相同，持续升温和持续降温却可能需要不同处理。序列网络、状态空间模型和时间核各有办法利用这种顺序。比较时继续检查它们如何表达依赖、需要多少计算，以及给出的不确定性究竟针对潜在规律还是未来读数。